% =====================================================================
% Conclusion - sources: all chapters; canonical figures per P0.
% =====================================================================

\section{Conclusion}
\label{sec:conclusion}

We presented the axiom-governance methodology of TOE-SYLVA, an
AI-assisted formal library that scaled past its own review
capacity, and the honesty layer that now disciplines it. The
methodology has four moving parts, each independently useful and
together closing the loop:

\begin{enumerate}
\item \textbf{A registry-first audit}: every \texttt{axiom}
declaration in the non-batch corpus is registered with a semantic
class (448 entries at baseline: primitive 250, definitional 142,
placeholder 42, schema 14), so that debt is explicit before it is
addressed.
\item \textbf{Redemption sweeps with semantic re-grading}: six
sweeps closed 196 entries ($30+44+31+29+30+32$), and the P0
reconciliation re-graded them into five classes --- vacuous 50,
$P\!\to\!P$ conditionalization 90, data-to-definition 17,
definitional reconstruction 24, substantive 13, plus 2 bundling
restructurings excluded from the content count. The honest
headline is 13.
\item \textbf{Three-caliber reconciliation}: registry (448),
live-declaration (253), and sweep-process (196) accounts now close
identically ($445-166-26=253$ and $448-253=195=142+42+11$), the
Agda side is pinned at 7 files / 25 blocks / 131 name-level
postulates, and the authoritative proof rate is 3.91\%
(22,328/571,611, handcrafted; batch reported separately).
\item \textbf{An honesty layer}: claim tiers with falsifiability
criteria, a caliber-discipline standard, a quarantine protocol,
and a CI gate that fails on silent deletions and missing markers.
\end{enumerate}

Two case studies showed the lifecycle end to end: a CLAIM$\to$
THEOREM upgrade along a Chern--Simons chain whose registered
theorem is deliberately labeled as what it is --- a
\texttt{mathlib}-delegated cast-summation identity, not
topological integrality --- and a Dedekind-domain segment that is
honest about contributing zero new mathematics while validating
the attribution mechanism.

For the growing population of AI-generated formal content, we
claim one general lesson. Every number this paper reports ---
including the unflattering ones --- is bound to a declared,
reproducible caliber. That binding, not any particular
percentage, is what makes a proof-rate metric trustworthy; and
the cheapest time to install it is before the corpus, and the
metrics, have a history worth laundering.
